% =============================================================================
% Capítulo 4 — Metodologia                                      (peso: pesado)
% -----------------------------------------------------------------------------
% Fonte principal: plano §6. Detalhes de protocolo entram à medida que os itens
% A07–A23 do context/action_log.md forem fechados (congelamento no mês 3).
% Taxonomia de validade: Wohlin et al. (R1).
% =============================================================================
\chapter{Metodologia}
\label{cap:metodologia}

O desenho experimental e a classificação das ameaças à validade seguem
\citeonline{wohlin2024}.
\pendencia{Visão geral: fase de solução seguida de dois benchmarks sobre as versões H e D.}

% -----------------------------------------------------------------------------
\section{Fase de solução}
\label{sec:fase-solucao}
% Fonte: plano §6.1. C4 + ADRs, harness como biblioteca,
% versão D com paridade funcional (A19), congelamento de versões.
\pendencia{Documentação arquitetural, implementação do harness e da \versaoD.}

\subsection{Documentação arquitetural}
\pendencia{C4 e ADRs — inclusive retry e logging.}

\subsection{Implementação do \harness}
\pendencia{MVP de 5 componentes; interface estável (A12).}

\subsection{Versão de controle (\versaoD) e paridade funcional}
\pendencia{Política da versão D (A07) e critério de paridade (A19).}

% -----------------------------------------------------------------------------
\section{Benchmark 1 --- Substituibilidade}
\label{sec:b1}
% Fonte: plano §6.2. Ancoragem: R14 (Mo), R13 (Kaur), R12 (Issarny).
% Pendências: A04 (provedores), A07, A08 (checklist), A09 (métrica de código),
% A10 (tempo).

\subsection{Variáveis e métricas}
\begin{table}[h!]
  \centering
  \Caption{\label{tab:b1-metricas}Métricas do Benchmark 1}
  \UFCtab{}{
    \begin{tabularx}{\textwidth}{lXl}
      \toprule
      Dimensão & Métrica & Papel \\
      \midrule \midrule
      Comportamento      & Casos de uso que continuam funcionando após a troca & Primária \\
      Alteração no código & Arquivos, linhas e módulos alterados               & Primária \\
      Impacto estrutural & Componentes ou interfaces afetados                  & Primária \\
      Dependência        & Referências diretas ao SDK do provedor              & Primária \\
      Testes             & Testes alterados ou quebrados                       & Primária \\
      Esforço            & Tempo de implementação da substituição              & Secundária \\
      \bottomrule
    \end{tabularx}
  }{
    \Fonte{Elaborado pelo autor}
  }
\end{table}

\subsection{Protocolo}
\pendencia{Troca de provedor em H e D; checklist pré-registrado (A08).}

\subsection{Mitigação do viés do experimentador}
\pendencia{Métricas objetivas como primárias; tempo secundário; ordem registrada.}

% -----------------------------------------------------------------------------
\section{Benchmark 2 --- Eficiência energética}
\label{sec:b2}
% Fonte: plano §6.3. Ancoragem: R2, R3, R4, R5, R6.
% Pendências: A11 (RAPL), A13 (determinismo), A14 (unidade amostral),
% A15–A18 (estatística), A21–A22 (piloto), A29 (modelo/serving).

\subsection{Matriz experimental}
\pendencia{2×2: requisição simples vs.\ agêntica × H vs.\ D.}

\subsection{Métricas e critério de aceitação}
\pendencia{J/requisição, J/token, potência, latência; gating de corretude.}

\subsection{Protocolo de telemetria e sessões}
\pendencia{NVML a 100 ms; baseline de ociosidade; estabilização térmica; intercalação H/D.}

\subsection{Análise estatística}
\pendencia{Pacote definitivo após A15–A18.}

% -----------------------------------------------------------------------------
\section{Ameaças à validade}
\label{sec:ameacas}
% Fonte: plano §6.4. Quatro blocos de \citeonline{wohlin2024}.
\subsection{Validade interna}
\pendencia{}
\subsection{Validade externa}
\pendencia{Inclui número de pontos de integração (A07).}
\subsection{Validade de construto}
\pendencia{}
\subsection{Validade de conclusão}
\pendencia{}

% -----------------------------------------------------------------------------
\section{Base metodológica reaproveitada}
\label{sec:base-reaproveitada}
% Fonte: plano §6.5. Citação de R7 padronizada após A06.
\pendencia{Protocolo de medição de trabalho prévio do autor; diferenças de infraestrutura.}
